\documentclass[12pt]{article}

\usepackage[T1]{fontenc}
\usepackage[utf8]{inputenc}
\usepackage[brazil]{babel}
\usepackage{lmodern}
\usepackage{microtype}
\usepackage[a4paper,margin=2.5cm]{geometry}
\usepackage{setspace}
\usepackage{indentfirst}

\setstretch{1.15}

\title{COMPREENDENDO CICLOS ECONÔMICOS \thanks{Texto Traduzido livremente por Luiz Mario Andrade com o auxílio da ferramenta Chat GPT 5.4. Artigo preparado para a Conferência de Kiel sobre Crescimento sem Inflação, 22-23 de junho de 1976; revisado em agosto de 1976. Gostaria de agradecer a Gary Becker, Jacob Frenkel, Don Patinkin, Thomas Sargent e Jose Scheinkman por seus comentários e sugestões.}}
\author{Robert E. Lucas, Jr. \\ Universidade de Chicago}
\date{Agosto de 1976}

\begin{document}

\maketitle

\section*{I.}

Por que ocorre que, em economias capitalistas, as variáveis agregadas sofrem flutuações repetidas em torno da tendência, todas essencialmente do mesmo caráter? Antes da \textit{Teoria Geral} de Keynes, a resolução dessa questão era considerada um dos principais desafios pendentes para a pesquisa econômica, e as tentativas de enfrentar esse desafio eram chamadas de teoria dos ciclos econômicos. Além disso, entre os teóricos dos ciclos econômicos do entre guerras, havia um amplo acordo sobre o que significaria resolver esse problema. Para citar Hayek como um exemplo proeminente:

\begin{quote}
[A] incorporação de fenômenos cíclicos no sistema da teoria do equilíbrio econômico, com a qual eles estão em "aparente contradição", permanece o problema crucial da Teoria do Ciclo Comercial;\footnote{Hayek (1933), p. 33n.}

Por "teoria do equilíbrio", entendemos aqui primordialmente a teoria moderna da interdependência geral de todas as quantidades econômicas, que foi expressa mais perfeitamente pela Escola de Lausanne de economia teórica.\footnote{Hayek (1933), p. 42n.}
\end{quote}

Uma consequência primária da Revolução Keynesiana foi o redirecionamento do esforço de pesquisa para longe dessa questão, voltando-se para a questão aparentemente mais simples da determinação do produto em um ponto no tempo, tomando a história como dada.\footnote{Este redirecionamento foi consciente e explícito por parte de Keynes. Veja, por exemplo, a primeira frase de seu capítulo sobre o ciclo econômico. "Visto que afirmamos ter demonstrado nos capítulos precedentes o que determina o volume de emprego em qualquer momento, segue-se que, se estivermos certos, nossa teoria deve ser capaz de explicar os fenômenos do ciclo econômico" (1936), p. 313.} Uma consequência secundária desta Revolução, devida mais a Tinbergen do que a Keynes, foi um rápido aumento no nível de precisão e explicitude com que as teorias econômicas agregadas eram formuladas. Como resultado, a macroeconomia keynesiana beneficiou-se de várias décadas de melhoria metodológica, enquanto, deste ponto de vista técnico, os esforços dos teóricos dos ciclos econômicos parecem irremediavelmente desatualizados.

No entanto, de outro ponto de vista, eles parecem bastante modernos. A observação de que a macroeconomia necessita de uma fundamentação microeconômica tornou-se comum e, embora haja muita confusão sobre a natureza dessa necessidade e sobre o que significaria satisfazê-la, é provável que muitos economistas modernos não tivessem dificuldade em aceitar a formulação do problema de Hayek como aproximadamente equivalente à sua própria. Quer isso seja verdade ou não, desejo argumentar neste ensaio que deveria ser assim, ou que o progresso mais rápido em direção a uma teoria econômica agregada coerente e útil resultará da aceitação da formulação do problema conforme avançada pelos teóricos dos ciclos econômicos, e não de novas tentativas de refinar as estruturas improvisadas às quais a macroeconomia keynesiana nos conduziu.

Honrar os ancestrais intelectuais é um objetivo válido em si mesmo, mas há uma razão mais imediata para interpretar a busca contemporânea por uma economia agregativa teoricamente sólida como uma retomada do trabalho dos teóricos pré-keynesianos. Acompanhando o redirecionamento do interesse científico ocasionado pela Revolução Keynesiana, houve uma mudança acentuada na natureza da contribuição para a política econômica que os economistas esperavam oferecer e que o público passou a aceitar amplamente. O esforço para "explicar os ciclos econômicos" tinha sido direcionado para identificar fontes institucionais de instabilidade, com a esperança de que, uma vez compreendidas, essas fontes pudessem ser removidas ou sua influência mitigada por mudanças institucionais apropriadas. O processo vislumbrado era o dolorosamente lento da discussão pública e reforma legislativa; por outro lado, havia a esperança de melhoria institucional de longo prazo ou "permanente". O abandono do esforço para explicar os ciclos econômicos acompanhou a crença de que a política econômica poderia efetuar movimentos imediatos, ou de curtíssimo prazo, da economia de um estado atual indesejável, independentemente de como se tenha chegado a ele, para um estado melhor.

A crença de que este último objetivo é alcançável, e que a tentativa de se aproximar de alcançá-lo é a única tarefa legítima da pesquisa em economia agregada, é tão difundida que o argumento em contrário é visto como "destrutivo", uma tentativa deliberada de tornar a vida mais difícil para os colegas que estão apenas tentando melhorar a sorte da humanidade. No entanto, a situação é simétrica. Se os teóricos dos ciclos econômicos estivessem corretos, a manipulação de curto prazo na qual grande parte da economia agregativa está agora focada apenas desvia a atenção da discussão de políticas de estabilização que poderiam realmente ser eficazes; tal adiamento é, além disso, acompanhado pela erosão constante e inteiramente compreensível na crença por parte dos não-economistas de que a economia agregativa tenha algo útil a dizer.

Na próxima seção, revisarei algumas das principais características qualitativas dos eventos que chamamos de ciclos econômicos e, em seguida, passarei para a resposta keynesiana a esses fatos, para o progresso feito ao longo da linha iniciada por Keynes e Tinbergen e, finalmente, para os severos limites a este progresso que agora se tornaram aparentes. O restante do ensaio considerará as perspectivas de explicar fenômenos cíclicos por uma teoria econômica, no sentido restrito em que Hayek e outros teóricos dos ciclos econômicos usaram esse termo.

\section*{II.}

Permitam-me começar a refinar a discussão revisando as principais características qualitativas das séries temporais econômicas que chamamos de "o ciclo econômico". Tecnicamente, os movimentos em torno da tendência no produto nacional bruto em qualquer país podem ser bem descritos por uma equação de diferença perturbada estocasticamente de ordem muito baixa. Esses movimentos não exibem uniformidade de período ou amplitude, o que quer dizer que não se assemelham aos movimentos ondulatórios determinísticos que às vezes surgem nas ciências naturais. Aquelas regularidades que são observadas residem nos comovimentos entre diferentes séries temporais agregativas.

As principais entre estas são as seguintes:\footnote{As características das séries temporais econômicas listadas aqui são, curiosamente, tanto "bem conhecidas" quanto caras para documentar de forma cuidadosa e abrangente. Uma introdução útil e substantivamente orientada é dada por Mitchell (1951), que resume principalmente a experiência dos EUA no entre guerras. A referência técnica básica para esses métodos é Burns e Mitchell (1946). A experiência monetária dos EUA é melhor exibida em Friedman e Schwartz (1963). Uma fonte inestimável para as séries britânicas anteriores é Gayer, Rostow e Schwartz (1953), esp. Vol. II. Os fenômenos documentados nessas fontes são, obviamente, muito mais amplamente observados. A maioria pode ser inferida, embora com alguma dificuldade, a partir da estrutura estimada de modelos econométricos modernos. Uma contribuição recente importante é Sargent e Sims (1976), que resume as séries trimestrais dos EUA do pós-guerra de várias maneiras sugestivas, levando a um quadro qualitativo muito próximo ao fornecido por Mitchell, mas dentro de um arcabouço estocástico explícito, de modo que seus resultados são replicáveis e criticáveis em um nível que os de Mitchell não são.} (i) Os movimentos do produto em setores amplamente definidos movem-se juntos. (Na terminologia de Mitchell, exibem alta conformidade; na linguagem moderna de séries temporais, têm alta coerência). (ii) A produção de bens duráveis de produção e consumo exibe muito mais amplitude do que a produção de não duráveis. (iii) A produção e os preços de bens agrícolas e recursos naturais têm conformidade inferior à média. (iv) Os lucros das empresas mostram alta conformidade e amplitude muito maior do que outras séries. (v) Os preços geralmente são procíclicos. (vi) As taxas de juros de curto prazo são procíclicas; as taxas de longo prazo o são ligeiramente. (vii) Os agregados monetários e as medidas de velocidade são procíclicos.

Não há, tanto quanto sei, necessidade de qualificar essas observações restringindo-as a países ou períodos de tempo específicos: elas parecem ser regularidades comuns a todas as economias de mercado descentralizadas. Embora não haja absolutamente nenhuma razão teórica para antecipar isso, os fatos levam a concluir que, com respeito ao comportamento qualitativo dos comovimentos entre séries, os ciclos econômicos são todos iguais. Para economistas com inclinação teórica, essa conclusão deve ser atraente e desafiadora, pois sugere a possibilidade de uma explicação unificada dos ciclos econômicos, fundamentada nas leis gerais que governam as economias de mercado, em vez de em características políticas ou institucionais específicas de países ou períodos particulares.

Omiti o comportamento das estatísticas de comércio exterior do catálogo acima de fenômenos a serem explicados, em parte porque, para uma grande economia como a dos EUA, as estatísticas comerciais não exibem conformidade alta o suficiente para serem ciclicamente interessantes. Para um país menor, certamente, os movimentos das exportações fariam muito para "explicar" os ciclos, mas focar em explicações de economia aberta seria, penso eu, fugir da questão mais difícil e crucial sobre as origens últimas dos movimentos cíclicos.

Também omitido, mas um fenômeno marcante demais para passar sem comentários, é a redução geral na amplitude de todas as séries nos vinte e cinco anos seguintes à Segunda Guerra Mundial. Nesse nível puramente descritivo, é impossível distinguir boa sorte de boa política. No entanto, um período tão longo de relativa estabilidade sugere fortemente que não há nada inerente ao funcionamento das economias de mercado que exija conviver com o nível de instabilidade que estamos vivenciando agora, ou ao qual estávamos sujeitos nos anos anteriores à Segunda Guerra Mundial. Ou seja, as tentativas de documentar e explicar movimentos cíclicos regulares não precisam estar conectadas de forma alguma a uma presunção de que tais movimentos sejam uma característica inevitável das economias capitalistas.

\section*{III.}

As implicações dos modelos macroeconômicos keynesianos conformam-se bem às características das séries temporais revisadas acima. Versões iniciais (por exemplo, por Hicks, 1937, e Modigliani, 1944) ajustam-se bem qualitativamente; os modelos econométricos que se desenvolveram a partir desta teoria e do trabalho inicial amplamente independente de Tinbergen\footnote{Por exemplo, veja Tinbergen (1939). Este trabalho não era explicitamente keynesiano; de fato, foi concebido como um complemento empírico à revisão e síntese de Haberler sobre o trabalho teórico sobre ciclos econômicos (1936). Keynes, por sua vez, foi ativamente hostil ao trabalho de Tinbergen. Veja Moggridge (1973), pp. 285-320. Ao referir-me àqueles que se basearam em parte no trabalho de Tinbergen como "keynesianos", estou, portanto, contribuindo para a continuação de uma injustiça histórica.} conformam-se bem quantitativamente. Esses modelos localizaram as perturbações primárias no comportamento do investimento, ligadas através de defasagens (no modelo dos EUA de Tinbergen) à série de lucros altamente volátil. Movimentos nessas séries de alta amplitude induzem então movimentos gerais no produto e no emprego. Como essas perturbações eram, nos termos de Hicks, "deslocamentos IS", elas eram consistentes com taxas de juros e velocidade movendo-se prociclicamente. A suposição de salários e preços rígidos foi uma boa primeira aproximação empírica. Mais tarde, um setor de salários-preços (mais tarde chamado de curva de Phillips) foi adicionado para ajustar os movimentos procíclicos observados de salários e preços.\footnote{Klein e Goldberger (1955).}

Nesta descrição, os movimentos na moeda não desempenham papel importante na explicação dos ciclos. Esta característica certamente não resultou diretamente dos modelos teóricos; Keynes, Hicks e Modigliani deram grande ênfase às forças monetárias. A desênfase na moeda deu-se por razões empíricas: os econometristas, a partir de Tinbergen, descobriram que os fatores monetários não pareciam muito importantes empiricamente.\footnote{Tinbergen (1939), pp. 183-185. Tinbergen, assim como a maioria dos macroeconometristas subsequentes, usou a significância das taxas de juros para testar a importância da moeda.}

O sucesso empírico desses desenvolvimentos foi medido de uma forma original e historicamente apta por Adelman e Adelman (1959) em sua simulação do modelo Klein-Goldberger da economia dos EUA. Os Adelmans colocaram, de forma precisa, a questão de se um observador armado com os métodos de Burns e Mitchell (1946) poderia distinguir entre uma coleção de séries econômicas geradas artificialmente por um computador programado para seguir as equações de Klein-Goldberger e as séries análogas geradas por uma economia real. A resposta, para a surpresa evidente dos Adelmans (e, suspeita-se, de Klein e Goldberger, que de forma alguma haviam direcionado seus esforços para atender a esse critério), foi não.\footnote{Não é correto que uma busca por "bons ajustes" teria levado a um modelo satisfazendo os critérios dos Adelmans; pense em ajustar polinômios no tempo para "explicar" cada série ao longo do período da amostra.}

Esta conquista sinalizou um novo padrão para o que significa entender os ciclos econômicos. Demonstra-se a compreensão dos ciclos econômicos construindo um \textit{modelo} no sentido mais literal: uma economia artificial plenamente articulada que se comporta ao longo do tempo de modo a imitar de perto o comportamento das séries temporais de economias reais. Os modelos macroeconômicos keynesianos foram os primeiros a atingir este nível de explicitude e precisão empírica; ao fazê-lo, alteraram o significado do termo "teoria" a tal ponto que as teorias de ciclos econômicos mais antigas não podiam realmente ser vistas como "teorias".

Esses modelos não são, no entanto, "teorias de equilíbrio" no sentido de Hayek. De fato, Keynes escolheu iniciar a \textit{Teoria Geral} com a declaração (o Capítulo II não é mais do que isso) de que uma teoria de equilíbrio era inalcançável: que o desemprego não era explicável como consequência de escolhas individuais e que a falha dos salários em se moverem como previsto pela teoria clássica deveria ser tratada como devida a forças além do poder da teoria econômica para iluminar.

Keynes escreveu como se a natureza "involuntária" do desemprego fosse verificável por observação direta, como se fosse possível olhar para um mercado e verificar diretamente se ele está em equilíbrio ou não. No entanto, havia razões empíricas sérias por trás dessa escolha, pois em lugar nenhum a "aparente contradição" entre "fenômenos cíclicos" e a teoria do "equilíbrio econômico" é mais aguda do que no comportamento do mercado de trabalho. Por que, diante de salários e preços nominais moderadamente flutuantes, as famílias escolheriam ofertar trabalho a taxas acentuadamente irregulares ao longo do tempo? A maioria dos teóricos dos ciclos econômicos evitou esse problema crucial, e aqueles que o abordaram não o resolveram. Keynes percebeu que, ao simplesmente contornar esse problema com o postulado não explicado de preços nominais rígidos, um modelo de outra forma clássico poderia ser transformado em um modelo que fazia um bom trabalho ao explicar as séries temporais observadas.

Essa decisão por parte do teórico mais prestigiado de sua época libertou uma geração de economistas da disciplina imposta pela teoria do equilíbrio e, como descrevi, essa liberdade foi rápida e frutuosamente explorada pelos macroeconometristas. Agora em posse de réplicas detalhadas e quantitativamente precisas da economia real, os economistas pareciam ter um meio barato de avaliar várias medidas de política econômica propostas. Parecia legítimo tratar as recomendações de política que emergiam deste procedimento como se tivessem sido testadas experimentalmente, mesmo que tais políticas nunca tivessem sido tentadas em nenhuma economia real.

No entanto, a capacidade de um modelo de imitar o comportamento real da maneira testada pelos Adelmans (1959) não tem quase nada a ver com sua capacidade de fazer previsões condicionais precisas, para responder a perguntas da forma: como o comportamento teria diferido se certas políticas tivessem sido diferentes em formas especificadas? Essa habilidade requer a invariância da estrutura do modelo sob variações de política do tipo que está sendo estudado. A invariância de parâmetros em um modelo econômico não é, obviamente, uma propriedade que possa ser assegurada antecipadamente, mas parece razoável esperar que nem os gostos nem a tecnologia variem sistematicamente com variações em políticas anticíclicas. Em contraste, as regras de decisão dos agentes mudarão, em geral, com mudanças no ambiente. Um modelo de equilíbrio é, por definição, construído de modo a prever como agentes com gostos e tecnologia estáveis escolherão responder a uma nova situação. Qualquer modelo de desequilíbrio, construído simplesmente codificando as regras de decisão que os agentes acharam útil utilizar ao longo de algum período amostral anterior, sem explicar por que essas regras foram usadas, não terá utilidade para prever as consequências de mudanças de política não triviais.

A importância quantitativa deste problema é, obviamente, uma questão a ser resolvida pelo exame de relações específicas em modelos específicos. Argumentei em outro lugar\footnote{Lucas (1976).} que ele é de importância fatal em virtualmente todos os setores dos modelos macroeconômicos modernos, primordialmente por causa do tratamento defeituoso das expectativas nesses modelos. Em vez de revisar esses argumentos em detalhes, permitam-me citar a ilustração mais gráfica: nossa experiência durante a recente "estagflação".

Até 1970, os principais modelos econométricos dos EUA implicavam que políticas monetárias e fiscais expansionistas que levassem a uma inflação sustentada de cerca de 4 por cento ao ano levariam também a taxas de desemprego sustentadas de menos de 4 por cento, ou cerca de um ponto percentual completo abaixo da média do desemprego durante qualquer longo período da história dos EUA.\footnote{Hirsch (1972), de Menil e Enzler (1972).} Essas previsões foram amplamente endossadas por muitos economistas que não estavam intimamente envolvidos com previsões econométricas. Anteriormente, Friedman (1968) e Phelps (1968) haviam argumentado, puramente com base na observação de que o comportamento de equilíbrio é invariante sob a mudança de unidades representada pela inflação sustentada, que nenhum decréscimo sustentado no desemprego resultaria da inflação sustentada. Nesse caso, o experimento de política em questão foi, muito infelizmente, realizado, e seu resultado é agora claro demais para exigir uma revisão detalhada.

É importante que a lição deste episódio não seja perdida. A questão é muito mais profunda do que a adição de algumas novas variáveis às curvas de Phillips econométricas (embora esta seja a única revisão nos modelos macroeconômicos que se seguiu a ela), como Friedman deixou claro em seu Discurso Presidencial. O argumento de Friedman não procedeu com base em um modelo agregativo específico, com um "setor salários-preços" melhor do que os modelos padrão. Pelo contrário, baseou-se em uma característica geral do equilíbrio econômico: a homogeneidade de grau zero das funções de demanda e oferta. Assim, sem usar nenhum modelo muito específico, e sem reivindicar a habilidade de prever detalhadamente a resposta inicial da economia a uma inflação, pode-se, no caso da inflação sustentada, raciocinar que, se a taxa de desemprego anterior à inflação fosse uma taxa de equilíbrio (ou "natural"), então a mesma taxa será um equilíbrio uma vez que a inflação esteja em curso.

O caso da inflação sustentada é relativamente simples (embora aparentemente não simples demais, pois ainda é altamente controverso). Para outros tipos de questões de política, precisar-se-ia de um modelo mais explícito. Como a variância, e outros momentos, do produto real mudariam se uma política de 4 por cento de crescimento monetário fosse adotada? Sob uma regra fiscal de orçamento equilibrado? Sob taxas de câmbio flexíveis em vez de fixas? Pode-se gerar respostas numéricas para questões desse tipo a partir dos modelos macroeconômicos atuais, mas não há razão para que alguém leve esses números a sério. Por outro lado, tampouco podem ser obtidas respostas quantitativas por raciocínio puramente teórico. Para obtê-las, é necessário um relato explícito e de equilíbrio do ciclo econômico.

\section*{IV.}

Resumi, na seção II, as principais características do comportamento cíclico em quantidades e preços. Na seção III, argumentei a necessidade prática de explicar esses fatos em termos de equilíbrio (isto é, não keynesianos). Ou seja, desejar-se-ia uma teoria que desse conta dos movimentos observados nas quantidades (emprego, consumo, investimento) como uma resposta otimizadora aos movimentos observados nos preços.

Na próxima seção, descreverei o ponto de vista geral em relação à tomada de decisão individual a ser adotado no restante do artigo e explicarei, em particular, por que o caráter recorrente dos ciclos econômicos é de importância central. Dado este ponto de vista geral, considerarei nas seções VI e VII a forma como os movimentos de preços relativos induzem flutuações no emprego e no investimento. As seções VIII, IX e X examinam as condições sob as quais essas mesmas respostas quantitativas podem ser desencadeadas por movimentos nos preços gerais, ou nominais. Sem surpresa, a fonte dos movimentos gerais de preços é localizada, na seção XI, em mudanças monetárias.

\section*{V.}

A visão do problema de decisão individual prototípico adotada pela teoria do capital moderna é um ponto de partida útil para considerar o comportamento ao longo do ciclo, embora seja em alguns aspectos altamente enganosa. Um agente inicia um período com estoques de vários tipos de capital acumulados no passado. Ele enfrenta trajetórias temporais de preços aos quais pode negociar no presente e no futuro. Com base em suas preferências sobre trajetórias temporais de trabalho ofertado e bens consumidos, ele formula um plano. Sob certeza, ele é visto como simplesmente executando um único plano sem revisão; com incerteza, ele deve elaborar um plano de contingência, dizendo como reagirá a eventos imprevistos.

Até para começar a pensar sobre problemas de decisão desta forma geral, é necessário imaginar uma visão razoavelmente precisa do futuro na mente desse agente. De onde ele obtém essa visão e como um observador pode inferir qual é ela? Esse aspecto do problema recebeu um tratamento um tanto descuidado na teoria do capital tradicional e nenhum tratamento na macroeconomia tradicional. Como isso é absolutamente crucial para a compreensão dos ciclos econômicos, devemos persegui-lo aqui em algum detalhe.

Em um nível puramente formal, sabemos que um agente racional deve formular uma distribuição de probabilidade conjunta subjetiva sobre todas as variáveis aleatórias desconhecidas que incidem sobre suas oportunidades de mercado presentes e futuras. A ligação entre esta visão subjetiva do futuro e a "realidade" é uma questão filosófica muito complexa, mas a forma como é resolvida tem pouco efeito sobre a estrutura do problema de decisão visto por um agente individual. Em particular, qualquer distinção entre tipos de aleatoriedade (como a distinção de Knight (1921) entre "risco" e "incerteza") é, nesse nível, sem sentido.

Infelizmente, a hipótese geral de que os agentes econômicos são tomadores de decisão bayesianos tem, em muitas aplicações, pouco conteúdo empírico: sem alguma forma de inferir qual é a visão subjetiva do futuro de um agente, essa hipótese não ajuda a entender seu comportamento. Até o comportamento psicótico pode ser (e hoje é) entendido como "racional", dada uma visão suficientemente anormal das probabilidades relevantes. Para praticar a economia, precisamos de alguma forma (curta de psicanálise, espera-se) de entender qual problema de decisão os agentes estão resolvendo.

John Muth (1961) propôs resolver esse problema identificando as probabilidades subjetivas dos agentes com as frequências observadas dos eventos a serem previstos, ou com probabilidades "verdadeiras", chamando a coincidência assumida de probabilidades subjetivas e "verdadeiras" de \textit{expectativas racionais}. Evidentemente, esta hipótese não terá valor para entender o comportamento psicótico. Nem será aplicável em situações em que não se pode adivinhar quais, se houver, frequências observáveis são relevantes: situações que Knight\footnote{Knight (1921). Estou interpretando a distinção risco-incerteza como se referindo não a uma classificação de diferentes tipos de problemas de decisão individual, mas à \textit{relação} entre o tomador de decisão e o observador.} chamou de "incerteza". Ela será muito provavelmente útil em situações em que as probabilidades de interesse dizem respeito a um evento recorrente razoavelmente bem definido, situações de "risco" na terminologia de Knight. Em situações de risco, a hipótese de comportamento racional por parte dos agentes terá conteúdo utilizável, de modo que o comportamento possa ser explicável em termos de teoria econômica. Em tais situações, as expectativas são racionais no sentido de Muth. Em casos de incerteza, o raciocínio econômico não terá valor.

Essas considerações explicam por que os teóricos dos ciclos econômicos enfatizaram o caráter recorrente do ciclo e por que devemos esperar que eles estivessem certos ao fazê-lo. Na medida em que os ciclos econômicos podem ser vistos como instâncias repetidas de eventos essencialmente semelhantes, será razoável tratar os agentes como reagindo às mudanças cíclicas como "risco", ou assumir que suas expectativas são racionais: que eles têm arranjos razoavelmente estáveis para coletar e processar informações, e que utilizam essa informação na previsão do futuro de uma forma estável, livre de vieses sistemáticos e facilmente corrigíveis.

\section*{VI.}

Ao passar dessas considerações gerais para uma teoria mais específica, será útil considerar como exemplo um agente "representativo".\footnote{Muitos dos argumentos nesta e em seções subsequentes foram desenvolvidos mais explicitamente em outros lugares. O tratamento paralelo mais próximo está em Lucas (1975). Veja também Phelps, et al. (1970), Barro (1976), Sargent e Wallace (1975), Sargent (1976). No que se segue, não documentarei argumentos particulares, nem tentarei repartir crédito (ou culpa) pelas ideias discutidas.} Imagine um único produtor-trabalhador, confrontado a cada período com um preço de mercado dado para um bem que ele então produz sob encomenda, a uma taxa fixa de produto por hora. Ou seja, ele chega ao seu local de trabalho, observa seu preço de venda atual, determina quantas horas trabalhar naquele dia, vende sua produção e então vai para casa relaxar.

O bem que ele recebe em troca do esforço é "dinheiro"; não me preocuparei com as razões históricas para esse arranjo, mas simplesmente o tomarei como dado. Este dinheiro, por sua vez, é gasto em uma ampla variedade de bens, diferentes de dia para dia. Algumas compras ele faz no caminho para casa, em um intervalo de uma hora do trabalho, ou vários dias depois. Assumo por enquanto que ele não detém outros títulos. Assumo também que este agente vive em um mundo livre de ciclos, no qual o nível geral ou médio de preços não muda, embora os preços individuais flutuem de dia para dia.

Agora, postulemos um aumento de 10 por cento no preço de venda de hoje, em comparação com a média dos preços passados. Como esse produtor hipotético responderá? A resposta dada pela teoria econômica deve ser: quem sabe? Neste ponto, não disse nada que permitisse imaginar o que o produtor pensa que este movimento de preço \textit{significa}. Se ele acredita que a mudança de preço sinaliza uma mudança permanente em seu preço de venda, sabemos por muitas evidências que ele não trabalhará mais e, provavelmente, trabalhará um pouco menos. Ou seja, sabemos que as elasticidades da oferta de trabalho de "longo prazo" (terminologia muito infeliz, já que a resposta de "longo prazo" a uma mudança permanente de preço será imediata) são zero ou negativas.

E se, no extremo oposto, a mudança de preço for transitória (como seria o caso se o preço de cada período fosse um sorteio independente de uma distribuição fixa)? A resposta neste caso equivale a conhecer a taxa na qual o produtor está disposto a substituir trabalho hoje por trabalho amanhã. Se o "lazer" for altamente substituível ao longo do tempo, ele trabalhará mais nos dias de preço alto e fechará cedo nos dias de preço baixo. Menos se sabe sobre as respostas reais da oferta de trabalho a movimentos transitórios de preços do que sobre a resposta de "longo prazo", mas o que sabemos indica que o lazer em um período é um excelente substituto para o lazer em outros períodos próximos. Evidências sistemáticas no nível agregado foram obtidas por Rapping e por mim (1970); Ghez e Becker (1975) chegaram à mesma conclusão em um nível desagregado. Os pequenos prêmios exigidos para induzir os trabalhadores a mudar feriados e férias (tirar a segunda-feira de folga em vez do sábado, duas semanas em março em vez de agosto) apontam para a mesma conclusão, e esta evidência "casual" é um pouco mais impressionante por causa de sua simplicidade probabilística: os feriados são conhecidos como transitórios. Com base nesta evidência, poder-se-ia prever uma resposta altamente elástica a mudanças transitórias de preços.

Antes de lidar com complicações a este exemplo, notemos sua promessa para a teoria dos ciclos econômicos. Descrevi um produtor que responde a pequenas flutuações de preços com grandes flutuações no produto e no emprego: exatamente o que observamos ao longo do ciclo. A descrição repousa sobre efeitos de substituição economicamente inteligíveis, não sobre "desequilíbrios" ininteligíveis. No entanto, vamos devagar: nossas observações agregativas referem-se a comovimentos de produto e preços \textit{geralmente}; o exemplo refere-se a movimentos de preços \textit{relativos} em um ambiente estacionário.

Antes de enfrentar essa difícil questão, consideremos algumas variações do exemplo acabado de considerar. Primeiro, de um ponto de vista descritivo, muitas vezes parece mais realista pensar na informação da demanda sendo transmitida aos produtores por mudanças nas \textit{quantidades}: novos pedidos, esgotamento de estoques e similares. Parece não haver nenhuma razão substantiva convincente para focar exclusivamente nos preços como sinais de demanda atual e futura. Nesse nível verbal, parece-me inofensivo e preciso usar os termos aumento de preço e aumento de vendas alternadamente. De forma um tanto surpreendente, no entanto, a análise rigorosa da determinação do equilíbrio quando os produtores estabelecem preços é extremamente difícil, e não existem exemplos relevantes para o comportamento do ciclo econômico.

Uma segunda variação é fácil de realizar. Em vez de considerar um trabalhador-empreendedor, poder-se-ia separar essas funções, introduzir firmas e considerar os mercados de trabalho e de produtos separadamente. No presente contexto, isso introduziria uma distinção entre salários e preços e levantaria a questão de arranjos de alocação de risco entre empregadores e trabalhadores.\footnote{Um desses arranjos é a prática de "demitir temporariamente" (\textit{laying off}) trabalhadores. Veja Azariadis (1975).} Também permitiria o estudo de conjuntos de informações possivelmente diferentes para firmas e trabalhadores. Nenhuma dessas questões é desprovida de interesse, mas todas são, em minha opinião, periféricas para a teoria do ciclo econômico. Os salários reais observados não são constantes ao longo do ciclo, mas também não exibem tendências consistentes pró ou contracíclicas. Isso sugere que qualquer tentativa de atribuir aos movimentos sistemáticos dos salários reais um papel central em uma explicação dos ciclos econômicos está condenada ao fracasso. Consequentemente, procederei como se o salário real fosse fixo, usando os termos "salários" e "preços" alternadamente.

Variações adicionais podem ser obtidas distinguindo entre vários usos do tempo do trabalhador-produtor quando ele não está trabalhando. Muitos autores tentaram, por exemplo, interpretar o desemprego medido como tempo dedicado à busca de emprego. Certamente, se alguém substitui o trabalho, substitui \textit{por} alguma outra atividade, e a experiência mostra que a crença de alguém na importância da substituição é reforçada por algumas ilustrações plausíveis. No entanto, há pouca evidência de que muito tempo seja gasto na busca de emprego, que a busca seja menos dispendiosa quando desempregado do que quando empregado ou, aliás, que o desemprego medido meça qualquer atividade sequer. Economicamente, a questão importante é a magnitude da elasticidade do emprego em relação a movimentos transitórios de salários e preços, não as razões pelas quais essa elasticidade é o que é.

De fato, suspeito que a indisposição de falar de trabalhadores em recessão como desfrutando de "lazer" é mais um testemunho da força da insistência de Keynes de que o desemprego é "involuntário" do que uma resposta a fenômenos observados. Não se quer sugerir que as pessoas \textit{gostam} de depressões! Obviamente, a hipótese de um mercado de trabalho em equilíbrio não carrega tal sugestão, não mais do que a observação de que as pessoas passam fome em mercados de alimentos em equilíbrio sugere que as pessoas gostam de passar fome.

\section*{VII.}

Variações mais complexas sobre este exemplo surgem quando o \textit{capital} de vários tipos é introduzido. Façamos isso, mantendo ainda a suposição de estabilidade ao longo do tempo no nível geral de preços.

Três possibilidades de interesse surgem. Primeiro, suponha que a produção atual possa ser armazenada como estoque de bens acabados. Essa possibilidade parece trabalhar \textit{contra} a explicação dos comovimentos preço-produto esboçada acima. O produtor certamente produzirá em períodos de preço baixo para venda posterior quando o preço estiver alto, suavizando a oferta de trabalho em relação ao caso em que o armazenamento é excluído. No nível industrial, no entanto, esse comportamento também amortece os movimentos de preços. O resultado líquido é provavelmente uma redução na elasticidade emprego-produção em relação ao preço e um aumento na elasticidade real vendas-preço.

Como uma segunda possibilidade, suponha que o produtor possa usar uma parte de sua produção atual para adquirir uma máquina que aumentará seu produto por hora em todos os períodos futuros. Como uma terceira, suponha que ele possa fazer um curso na escola que terá o mesmo efeito. Como essas duas possibilidades não diferem economicamente, elas podem ser consideradas como uma só. No exemplo de movimentos de preços puramente transitórios, discutido anteriormente, é claro que nenhuma dessas opções será exercida — desde que o produtor estivesse satisfeito com seu estoque original de capital. No momento em que o novo capital puder ser aplicado à produção, o movimento de preços que o fez parecer lucrativo terá desaparecido.

Movimentos de preços relativos atuais terão seu efeito máximo na acumulação de capital quando, no extremo oposto, forem considerados permanentes. Neste caso, no entanto, como observei, o emprego será insensível aos movimentos de preços. Assim, para observar o investimento e o emprego movendo-se sistematicamente na direção dos movimentos de preços relativos, deve ocorrer que tais movimentos sejam uma mistura de elementos transitórios e permanentes. Em tal situação, o produtor se verá obrigado a se engajar no que os engenheiros chamam de "processamento de sinais": ele observa uma única variável (preço) mudando ao longo do tempo; esses movimentos surgem de movimentos em variáveis mais fundamentais (os componentes transitórios e permanentes do preço) que não podem ser observados diretamente; a partir desses movimentos de preços observados, juntamente com seu conhecimento da importância relativa das duas fontes não observadas de mudança de preço, ele infere imperfeitamente os movimentos nos dois componentes. Com base em sua solução para este cálculo implícito de probabilidade condicional, ele toma uma decisão. Não surpreendentemente, a decisão acaba sendo uma média das decisões apropriadas para os dois extremos.

Para recapitular, assume-se que nosso produtor hipotético enfrenta uma variabilidade estocástica de preços, que é descritível como uma mistura de componentes transitórios e permanentes, ambos não observados. Sua resposta ótima aos movimentos de preços depende de dois fatores: a maneira como ele interpreta a informação contida nessas mudanças e suas preferências relativas à substituição intertemporal de lazer e consumo. Sob suposições consistentes com o comportamento racional e a evidência disponível, sua resposta a um aumento imprevisto de preço é um aumento considerável no trabalho ofertado, um declínio no estoque de bens acabados e uma expansão na acumulação de capital produtivo de todos os tipos. Esse comportamento é simétrico; as respostas às quedas de preços são o oposto.\footnote{O que está acontecendo com os gastos de consumo à medida que essas respostas de emprego e investimento ocorrem? Em sua crítica aos modelos de ciclo econômico de equilíbrio, Grossman (1973) argumenta que o consumo deve necessariamente mover-se na direção oposta ao trabalho ofertado. Como não é isso que se observa de fato ao longo do ciclo, seria de fato um paradoxo sério se uma correlação negativa fosse uma consequência da teoria da utilidade. Pode-se derivá-la para casos especiais (ver Lucas, 1972, Fig. 1), mas essa implicação certamente não é um fato geral para famílias otimizadoras; não decorre, por exemplo, da teoria de Rapping e minha (1970) ou da de Ghez e Becker (1975, cap. 4).}

\section*{VIII.}

É hora de pensar em situar este produtor representativo em uma economia composta por agentes semelhantes, embora, é claro, produzindo bens diferentes e sujeitos a diferentes movimentos de preços individuais. Para fazer isso, deve-se ir além dos movimentos de preços para as mudanças na tecnologia e no gosto que os fundamentam. Essas mudanças estão ocorrendo o tempo todo e, de fato, sua importância para os agentes individuais domina em muito os movimentos relativamente menores que constituem o ciclo econômico. No entanto, esses movimentos deveriam, em geral, levar a movimentos de preços relativos, não gerais. Uma nova tecnologia, reduzindo os custos de produção de um bem antigo ou tornando possível a produção de um novo, atrairá recursos para o bem beneficiado e os afastará da produção de outros bens. Mudanças de gosto em favor da compra de um bem envolvem gastos reduzidos em outros. Além disso, em uma economia moderna complexa, haverá um grande número de tais deslocamentos em qualquer período dado, cada um pequeno em importância em relação ao produto total. Haverá muita "compensação média" (\textit{averaging out}) de tais efeitos entre os mercados.

O cancelamento deste tipo é, penso eu, a razão mais importante pela qual não se pode buscar uma explicação dos movimentos gerais que chamamos de ciclos econômicos na mera presença, \textit{per se}, da imprevisibilidade das condições nos mercados individuais. No entanto, este argumento não é inteiramente rígido. É certamente possível que um grande número de agentes sinta espontaneamente um desejo de aumentar suas semanas de trabalho e expandir os investimentos. Mais seriamente, houve muitos casos de choques de oferta que afetam todos, ou muitos, setores da economia simultaneamente. Tais choques não se cancelarão da maneira que descrevi e induzirão flutuações de produto no agregado. Eles não levarão, no entanto, a movimentos que se ajustem à descrição esboçada na seção II: todos os deslocamentos de oferta levarão a movimentos de preços contracíclicos (mantidas as demais condições), em contraste com os movimentos procíclicos que observamos.

É, então, possível situar nosso produtor hipotético em um cenário de equilíbrio geral, no qual seu preço e produto flutuam, mas os níveis agregados não. Suas respostas a esses movimentos de preços relativos imitarão as respostas agregadas aos movimentos de preços gerais que constituem o ciclo econômico. Temos então um modelo coerente, mas não um que ainda dê conta dos fenômenos gerais a serem explicados. Este modelo pode, sem dificuldade, ser modificado para permitir flutuações de produto agregadas induzidas pela oferta, mas estas não guardam semelhança com o ciclo econômico moderno.

Antes de deixar este mundo de agregados estáveis, vale a pena enfatizar que a maior parte do risco que preocupa e desafia os agentes econômicos estaria presente em tal cenário. Os consumidores aceitarão um novo design de automóvel ou ele se tornará uma piada nacional? Uma dúzia de anos de treinamento em piano levarão ao palco de concertos ou apenas a um hobby prazeroso? Os salários de horas extras desta semana ajudarão a financiar a educação de um filho ou sustentarão a família durante a greve do mês que vem? No momento em que se adquire a informação necessária para resolver questões como essas, é tarde demais; de um jeito ou de outro, a pessoa já está comprometida.

Comparado a riscos desta natureza e magnitude, a questão de se as horas realmente trabalhadas no ano à frente serão 1,03 vezes o que se planeja por agora, ou 0,97, parece menor, e parece assim porque de fato é. Na teoria econômica agregativa, estamos acostumados a pensar nos ciclos econômicos como um tipo de risco imposto a um ambiente de outra forma estável. Tais hábitos de pensamento refletem a transferência de abstrações úteis para alguns propósitos para contextos onde envolvem distorções fatais da realidade.

\section*{IX.}

Abandonemos agora a suposição de estabilidade nos preços médios. Do ponto de vista do produtor individual, isso envolve apenas uma ligeira mudança na natureza do problema de processamento de sinais que deve ser resolvido. Antes, um dado movimento em seu "próprio preço" poderia significar uma mudança permanente de preço relativo ou uma transitória. Agora, também pode significar que \textit{todos} os preços estão mudando, uma situação que, se diagnosticada corretamente, não levaria a nenhuma resposta real por parte do produtor. No entanto, pela mesma razão que movimentos permanentes e transitórios de preços relativos não podem ser selecionados com certeza no momento, tampouco podem movimentos relativos e gerais ser distinguidos. Aumentos gerais de preços, exatamente como aumentos de preços relativos, induzirão movimentos na mesma direção no emprego e no investimento.

Diferentemente das respostas às mudanças de gosto e tecnologia descritas anteriormente, essas respostas aos aumentos gerais de preços não tenderão a se cancelar entre os mercados. Com certeza, alguns produtores observarão declínios na demanda mesmo durante expansões de preços, mas mais observarão aumentos (isso é o que um aumento geral de preços significa) e, portanto, mais estarão expandindo em termos reais do que contraindo. O efeito líquido será de comovimentos em preços, produto e investimento no nível agregado, exatamente como observado no ciclo real.

É essencial para este argumento que os movimentos gerais de preços não sejam percebidos como tais enquanto estão ocorrendo. No contexto dos modelos agregativos ordinariamente usados, essa suposição pode parecer implausível: como poderiam os negociantes não saber o preço dos bens? Na realidade de um mundo multi-mercadorias, no entanto, ninguém desejaria observar todos os preços todos os dias, nem muitos negociantes achariam os índices de preços publicados particularmente úteis. Um negociante otimizador processará aqueles preços de maior importância para seu problema de decisão com mais frequência e cuidado, aqueles de menor importância menos, e a maioria dos preços nem um pouco. Das muitas fontes de risco de importância para ele, o ciclo econômico e o comportamento agregado em geral não são, para a maioria dos agentes, de importância especial, e não há razão para que os negociantes especializem seus próprios sistemas de informação para diagnosticar movimentos gerais corretamente.

Pelo mesmo raciocínio, pode-se ver que a inflação sustentada não afetará as decisões reais dos agentes da mesma forma que os movimentos de preços transitórios o fazem. Nada é mais fácil do que detectar e corrigir o viés sistemático nas previsões. Tais correções não envolvem mudanças nos sistemas de informação dos agentes ou nos custos de processamento de informações. Pode haver, é claro, algum atraso no diagnóstico da inflação sustentada pelo que ela é; com a mesma frequência, os agentes perceberão incorretamente uma inflação transitória como se fosse sustentada.

Mudanças no grau de variabilidade dos preços terão efeitos mais fundamentais no comportamento de processamento de informações dos agentes, porque afetam os "pesos" colocados na informação de preços ao prever preços futuros. A ideia geral é que se confia menos em sinais de preços "ruidosos".

\section*{X.}

A resposta agregada ou média aos movimentos gerais de preços torna-se mais complexa à medida que se considera o investimento, bem como as respostas do emprego. As decisões de investimento serão distorcidas por movimentos gerais de preços, pelas mesmas razões que o emprego o será, e na mesma direção que as respostas induzidas por movimentos de preços relativos.

Complicações adicionais seguem-se, no entanto, da observação de que o investimento atual afeta a capacidade futura e, portanto, os preços futuros. Este efeito pode ser visto como uma extensão no tempo, talvez até uma amplificação, dos efeitos iniciais dos movimentos gerais de preços.

Para detalhar isso, imagine que algum evento ocorra que, se percebido corretamente por todos, induziria um aumento nos preços de forma geral. Cedo ou tarde, então, esse ajuste ocorrerá. Inicialmente, no entanto, mais negociantes do que o contrário percebem um movimento de preço relativo, possivelmente permanente, a seu favor. Como resultado, o emprego e o investimento aumentam. Com o tempo, à medida que a informação de preços se difunde pela economia, esses negociantes verão que se enganaram. Nesse meio tempo, no entanto, a capacidade adicionada retarda os aumentos de preços de forma geral, adiando o reconhecimento do choque inicial. Desta forma, choques assistemáticos ou de curto prazo nos preços podem levar a oscilações muito mais longas nos preços.

Além disso, há uma desaceleração automaticamente embutida nessa expansão de capacidade. Quando o reconhecimento da inflação geral ocorre, o investimento terá que se tornar menor que o normal por um tempo enquanto a capacidade se reajusta para baixo. Não há razão para esperar que esse reajuste venha rapidamente, ou que seja descritível como um "crash" ou "quebra" (\textit{bust}).

Este cenário, como a descrição anterior da resposta do emprego, depende crucialmente da confusão por parte dos agentes entre movimentos de preços relativos e gerais. Isso é especialmente claro no caso do investimento, já que a política de investimento ótima tem uma grande quantidade de "suavização" embutida nela: como o investimento é um compromisso de longo prazo, ele responderá apenas ao que parecem ser mudanças de preços relativos relativamente permanentes.

Esta observação levou, por razões sérias, ao ceticismo quanto à importância dos efeitos aceleradores no ciclo econômico. Como podem movimentos cíclicos moderados nos preços levar aos movimentos de alta amplitude observados nas compras de bens duráveis? Aqui, novamente, deve-se insistir na pequena contribuição do risco em toda a economia para a situação de risco geral enfrentada pelos agentes. Para projetos de investimento individuais, as taxas de retorno são altamente variáveis, muitas vezes negativas e frequentemente medidas em centenas por cento. Uma resposta rápida e atual ao que parece a outros um sinal "fraco" é muitas vezes a chave para um investimento bem-sucedido. O agente que espera até que a situação esteja clara para todos chega tarde demais; outra pessoa já adicionou a capacidade para atender à alta demanda. O que parece, no nível agregado, ser um padrão de resposta de alta amplitude a choques de baixa amplitude é, no nível em que as decisões são tomadas, uma resposta de alta amplitude a movimentos de amplitude ainda maior nos retornos de investimentos individuais.\footnote{A teoria do ciclo econômico "austríaca" ou de "excesso de investimento monetário" (veja Haberler, 1936, ou Hayek, 1933) baseava-se nessa mesma ideia de decisões de investimento equivocadas desencadeadas por sinais de preços espúrios. No entanto, o preço que essa teoria enfatizava era a taxa de juros, em vez dos preços dos produtos conforme enfatizado aqui. Dada a amplitude cíclica das taxas de juros, a elasticidade investimento-juros necessária para explicar a amplitude observada no investimento é alta demais para ser consistente com outras evidências.}

\section*{XI.}

Comecei a seção II com uma definição de ciclos econômicos como flutuações repetidas no emprego, produto e composição do produto, associadas a um certo padrão típico de comovimentos em preços e outras variáveis. Como em uma economia competitiva o emprego e o produto de vários tipos são escolhidos pelos agentes em resposta aos movimentos de preços, pareceu apropriado começar racionalizando os movimentos de quantidade observados como respostas racionais ou ótimas aos movimentos de preços observados. Isso foi realizado nas cinco seções anteriores. Volto-me agora para as fontes dos movimentos de preços.

Para explicar os movimentos seculares nos preços de forma geral, os movimentos seculares na quantidade de moeda funcionam extremamente bem. Este fato é tão bem estabelecido quanto qualquer outro que conhecemos na economia agregativa e não é sensível à forma como se mede os preços ou a quantidade de moeda.\footnote{Friedman e Schwartz (1963).} Não há dúvida séria quanto à direção do efeito nesta relação; ninguém argumenta que a antecipação da inflação do século XVI enviou Colombo ao Novo Mundo para localizar o ouro para financiá-la. Esta evidência não tem conexão direta com os ciclos econômicos, já que se refere a médias sobre períodos muito mais longos, mas as conexões indiretas são fortes demais para serem ignoradas: explicamos o padrão de comovimentos entre variáveis reais ao longo do ciclo como respostas a movimentos gerais de preços; sabemos que, no "longo prazo", os movimentos gerais de preços surgem primordialmente de mudanças na quantidade de moeda. Além disso, os movimentos cíclicos na moeda são grandes o suficiente para serem quantitativamente interessantes. Todos esses argumentos apontam para um choque monetário como a força que desencadeia o ciclo econômico real.

A evidência direta sobre correlações de curto prazo entre moeda, produto e preços é muito mais difícil de interpretar. Certos episódios extremos parecem indicar que depressões e recuperações são induzidas pela moeda.\footnote{Novamente, veja Friedman e Schwartz (1963).} Em geral, no entanto, concorda-se que o vínculo entre a moeda e estas e outras variáveis está sujeito, nos termos de Friedman, a "defasagens longas e variáveis".

Paradoxalmente, essa fraqueza na evidência de curto prazo ligando a moeda à atividade econômica, e em particular aos preços, é encorajadora do ponto de vista da teoria monetária do ciclo econômico. Para ver por que, recorde a ligação teórica entre movimentos gerais de preços e atividade econômica conforme esboçado acima. Essa conexão repousava na hipótese de que o problema de processamento de sinais de identificar movimentos gerais de preços a partir de observações de alguns preços individuais era \textit{difícil demais} para ser resolvido perfeitamente pelos agentes. Agora, suponha que fosse verdade que se pudesse descrever movimentos gerais de preços de curto prazo por uma função simples e fixa de movimentos defasados em algum agregado monetário publicado. Então, longe de ser difícil, o problema de processamento de sinais a ser resolvido pelos agentes seria \textit{trivial}; eles poderiam simplesmente observar os agregados monetários atuais, calcular os movimentos de preços atuais e futuros previstos que eles implicam e corrigir seu comportamento perfeitamente para essas mudanças de unidade. O resultado seria uma relação muito estreita entre moeda e preços, mesmo em períodos muito curtos, e nenhuma relação entre esses movimentos e mudanças nas variáveis reais.

Essas observações não explicam, é claro, \textit{por que} os efeitos monetários funcionam com defasagens longas e variáveis. Sobre essa questão pouco se sabe. Parece provável que a resposta resida na observação de que uma expansão monetária pode ocorrer de várias formas, dependendo da maneira como a moeda é "injetada" no sistema, com diferentes implicações de resposta de preços dependendo de qual via é selecionada. Isso sugeriria que se deveria descrever o "estado" monetário da economia como sendo determinado por algum agregado monetário não observável, fracamente relacionado aos agregados observados em períodos curtos, mas intimamente relacionado de forma secular.

\section*{XII.}

Permitam-me recapitular as principais características da teoria do ciclo econômico esboçada nas seções precedentes. Começamos imaginando uma economia com gostos e tecnologia flutuantes, implicando preços relativos em mudança contínua, e estudamos os comovimentos em quantidades e preços que surgiriam se os agentes se comportassem em seu próprio interesse e utilizassem suas informações incompletas de forma eficaz. Em seguida, sobrepusemos a essa economia movimentos assistemáticos e consideráveis em um agregado monetário, adicionando uma fonte adicional de "ruído" aos movimentos de preços individuais. O resultado é gerar um padrão de comovimentos entre séries agregadas que parece corresponder às observações resumidas na seção II.

Em retrospectiva, esse relato parece um tanto embaraçosamente simples: alguém se pergunta por que parece ser necessário desfazer uma Revolução para chegar a ele. No entanto, deve-se ter cuidado para não exagerar o que foi, de fato, alcançado. Penso que está bastante claro que não há nada no comportamento das séries temporais econômicas observadas que impeça ordená-las em termos de equilíbrio, e exemplos teóricos suficientes existem para dar confiança na esperança de que isso possa ser feito de forma explícita e rigorosa. Até o momento, no entanto, nenhum modelo de equilíbrio foi desenvolvido que atenda a esses padrões e que, ao mesmo tempo, pudesse passar no teste proposto pelos Adelmans (1959). Meu palpite seria que o sucesso nesse sentido está a cinco, mas não a vinte e cinco anos de distância.\footnote{Avançando mais neste palpite, é provável que tal modelo "bem-sucedido" seja um descendente próximo do de Sargent (1976).}

As implicações para a política econômica de uma teoria do ciclo econômico bem-sucedida do tipo aqui delineado são, creio, fáceis de adivinhar, mesmo quando a própria teoria está em um estágio preliminar. De fato, muito do que foi exposto acima é simplesmente uma tentativa de compreender e tornar mais explícito o modelo implícito subjacente às propostas de política de Henry Simons, Milton Friedman e outros críticos da política agregativa ativista. Ao buscar um relato de equilíbrio para os ciclos econômicos, aceitam-se antecipadamente limitações severas ao escopo da política governamental anticíclica que poderia ser racionalizada pela teoria. Na medida em que as flutuações são induzidas por instabilidade monetária gratuita, que não serve a nenhum propósito social, então o aumento da estabilidade monetária promete reduzir a variabilidade agregada real e aumentar o bem-estar. Não há dúvida, no entanto, de que \textit{alguma} variabilidade real permaneceria mesmo sob as políticas monetárias e fiscais mais suaves. Não há um caso \textit{prima facie} de que essa variabilidade residual seria melhor tratada por políticas governamentais centralizadas do que por respostas individuais descentralizadas.\footnote{Ou seja, a política anticíclica ativa exigiria o mesmo tipo de defesa de custo-benefício usada na avaliação de outros tipos de políticas governamentais. Veja Phelps (1972) e também a revisão de Prescott (1975).}

Diante dessa falta de novidade no campo da política, parece uma questão justa perguntar: por que precisamos da teoria? A resposta geral, penso eu, é que em uma sociedade democrática não basta acreditar que se está certo; deve-se ser capaz de explicar por que se está certo. Vivemos em uma sociedade na qual a taxa de desemprego flutua entre, digamos, 3 e 10 por cento. Segue-se que ambas as situações são alcançáveis, e é claro que a maioria das pessoas é mais feliz aos três do que aos dez. Também está claro que as políticas governamentais têm muito a ver com qual dessas situações prevalece em qualquer momento particular. O que poderia ser mais natural, então, do que ver a tarefa da economia agregativa como a de descobrir quais políticas levarão à situação mais desejável e, então, advogar sua adoção? Esta era a promessa da economia keynesiana e, mesmo agora, quando o esvaziamento científico desta promessa é mais evidente, seu apelo é compreensível para todos que compartilham a esperança de que as ciências sociais ofereçam mais do que uma elegante racionalização do estado de coisas existente.

O público economicamente alfabetizado teve cerca de quarenta anos para se sentir confortável com duas ideias relacionadas: que as economias de mercado são inerentemente sujeitas a flutuações violentas que só podem ser eliminadas por respostas governamentais flexíveis e enérgicas; e que os economistas estão em posse de um corpo de conhecimento cientificamente testado que lhes permite determinar, a qualquer momento, quais devem ser essas respostas. É duvidoso que muitos que não estejam profissionalmente comprometidos mantenham hoje a segunda dessas crenças. Isso por si só pouco resolve na disputa sobre se o papel do governo na política de estabilização deveria ser o de reduzir sua própria parte disruptiva ou compensar ativamente a instabilidade do setor privado. Enquanto o ciclo econômico permanecer "em aparente contradição" com a teoria econômica, ambas as posições parecem sustentáveis. Não parece haver forma de determinar como os ciclos econômicos devem ser tratados sem entender o que são e como ocorrem.

\section*{REFERÊNCIAS}

\begin{enumerate}
    \item Adelman, I., e Adelman, F. L., "The Dynamic Properties of the Klein-Goldberger Model." \textit{Econometrica}, 27, No. 4, (Outubro 1959), 596-625.
    \item Azariadis, C., "Implicit Contracts and Underemployment Equilibria," \textit{Journal of Political Economy}, 83, No. 6, (Dezembro 1975), 1183-1202.
    \item Barro, R. J., "Rational Expectations and the Role of Monetary Policy," \textit{Journal of Monetary Economics}, 2, No. 1, (Janeiro 1976), 1-32.
    \item Burns, A. F., e Mitchell, W. C. \textit{Measuring Business Cycles}. New York: National Bureau of Economic Research, 1946.
    \item Friedman, M., "The Role of Monetary Policy," Discurso Presidencial para a American Economic Association, \textit{American Economic Review}, 58, No. 1, (Março 1968), 1-17.
    \item Friedman, M., e Schwartz, A. J. \textit{A Monetary History of the United States, 1867-1960}. Princeton: Princeton University Press para o National Bureau of Economic Research, 1963.
    \item Gayer, A. D., Rostow, W. W., e Schwartz, A. J. \textit{The Growth and Fluctuation of the British Economy, 1790-1850}. Oxford: The Clarendon Press, 1953.
    \item Ghez, G. R., e Becker, G. S. \textit{The Allocation of Time and Goods Over the Life Cycle}. New York: National Bureau of Economic Research, 1975.
    \item Grossman, H. I., "Aggregate Demand, Job Search, and Employment," \textit{Journal of Political Economy}, 81, No. 6, (Novembro/Dezembro 1973), 1353-1369.
    \item Haberler, G. \textit{Prosperity and Depression}. Geneva: League of Nations, 1936.
    \item Hayek, F. A. von. \textit{Monetary Theory and the Trade Cycle}. London: Jonathan Cape, 1933.
    \item Hicks, J. R., "Mr. Keynes and the 'Classics': A Suggested Interpretation," \textit{Econometrica}, 5, (1937), 147-159.
    \item Hirsch, A. A., "Price Simulations with the OBE Econometric Model." in \textit{The Econometrics of Price Determination Conference}, (ed. O. Eckstein), Washington, D. C.: Board of Governors of the Federal Reserve System and Social Science Council, 1972.
    \item Keynes, J. M. \textit{The General Theory of Employment, Interest and Money}. London: Macmillan, 1936.
    \item Klein, L. A., e Goldberger, A. S. \textit{An Econometric Model of the United States, 1929-52}. Amsterdam: North Holland, 1955.
    \item Knight, F. H. \textit{Risk, Uncertainty and Profit}. Boston: Houghton Mifflin, 1921.
    \item Lucas, R. E., Jr., "Expectations and the Neutrality of Money," \textit{Journal of Economic Theory}, 4, No. 2, (Abril 1972), 103-123.
    \item \_\_\_\_\_\_\_\_\_\_, "An Equilibrium Model of the Business Cycle," \textit{Journal of Political Economy}, 83, No. 6, (Dezembro 1975), 1113-1144.
    \item \_\_\_\_\_\_\_\_\_\_, "Econometric Policy Evaluations: A Critique," in \textit{The Phillips Curve and Labor Markets}, (eds. K. Brunner e A. H. Meltzer), Carnegie-Rochester Conference Series on Public Policy, 1, Amsterdam: North Holland, 1976, 19-46.
    \item Lucas, R. E., Jr., e Rapping, L. A., "Real Wages, Employment, and Inflation," in \textit{Microeconomic Foundations of Employment and Inflation Theory}, (eds. E. S. Phelps, et al.), New York: Norton, 1970.
    \item de Menil, G., e Enzler, J. J., "Price and Wages in the FR-MIT Econometric Model," in \textit{The Econometrics of Price Determination Conference}, (ed. O. Eckstein), Washington, D. C.: Board of Governors of the Federal Reserve System and Social Science Council, 1972.
    \item Mitchell, W. C. \textit{What Happens During Business Cycles}. New York: National Bureau of Economic Research, 1951.
    \item Modigliani, F., "Liquidity Preference and the Theory of Interest and Money," \textit{Econometrica}, 12, No. 1, (Janeiro 1944), 45-88.
    \item Moggridge, D. (ed.) \textit{The Collected Writings of John Maynard Keynes}, Vol. XIV. London: Macmillan, 1973.
    \item Muth, J., "Rational Expectations and the Theory of Price Movements," \textit{Econometrica}, 29, No. 3, (Julho 1961), 315-335.
    \item Phelps, E. S., "Money Wage Dynamics and Labor Market Equilibrium," \textit{Journal of Political Economy}, 76, No. 4, II, (Julho/Agosto 1968), 687-711.
    \item \_\_\_\_\_\_\_\_\_\_. \textit{Inflation Policy and Unemployment Theory: The Cost-Benefit Approach to Monetary Planning}. London: Macmillan, 1972.
    \item Phelps, E. S., et al. \textit{Microeconomic Foundations of Employment and Inflation Theory}. New York: Norton, 1970.
    \item Prescott, E. C., "Efficiency of the Natural Rate," \textit{Journal of Political Economy}, 83, No. 6, (Dezembro 1975), 1229-1236.
    \item Sargent, T. J., "A Classical Macroeconometric Model for the United States," \textit{Journal of Political Economy}, 84, No. 2, (Abril 1976), 207-237.
    \item Sargent, T. J., e Sims, C. A., "Business Cycle Modeling Without Pretending to Have Too Much A Priori Economic Theory," University of Michigan working paper, Março 1976.
    \item Sargent, T. J., e Wallace, N., "'Rational' Expectations, the Optimal Monetary Instrument, and the Optimal Money Supply Rule," \textit{Journal of Political Economy}, 83, No. 2, (Abril 1975), 241-254.
    \item Tinbergen, J. \textit{Business Cycles in the United States of America, 1919-32}. Geneva: League of Nations, 1939.
\end{enumerate}

\end{document}